\documentclass[12pt,a4paper]{article}

\usepackage[utf8]{inputenc}
\usepackage[spanish]{babel}
\usepackage[T1]{fontenc}
\usepackage{geometry}
\usepackage{listings}
\usepackage{xcolor}
\usepackage{hyperref}
\usepackage{booktabs}
\usepackage{array}
\usepackage{enumitem}
\usepackage{fancyhdr}
\usepackage{mdframed}
\usepackage{tcolorbox}
\usepackage{amsmath}

\geometry{margin=2.5cm}

% ── Colores ────────────────────────────────────────────────────────────────────
\definecolor{codebg}{HTML}{F5F5F5}
\definecolor{codeframe}{HTML}{CCCCCC}
\definecolor{keywordcolor}{HTML}{0000CC}
\definecolor{stringcolor}{HTML}{008000}
\definecolor{commentcolor}{HTML}{888888}
\definecolor{alertbg}{HTML}{FFF8E1}
\definecolor{alertborder}{HTML}{F9A825}
\definecolor{taskbg}{HTML}{E8F5E9}
\definecolor{taskborder}{HTML}{2E7D32}

% ── Estilos de código ──────────────────────────────────────────────────────────
\lstdefinestyle{python}{
  language=Python,
  backgroundcolor=\color{codebg},
  frame=single,
  rulecolor=\color{codeframe},
  basicstyle=\ttfamily\small,
  keywordstyle=\color{keywordcolor}\bfseries,
  stringstyle=\color{stringcolor},
  commentstyle=\color{commentcolor}\itshape,
  breaklines=true,
  showstringspaces=false,
  numbers=left,
  numberstyle=\tiny\color{gray},
  numbersep=8pt,
  tabsize=4
}

\lstdefinestyle{texto}{
  backgroundcolor=\color{codebg},
  frame=single,
  rulecolor=\color{codeframe},
  basicstyle=\ttfamily\small,
  breaklines=true,
  showstringspaces=false
}

% ── Cajas de consigna ──────────────────────────────────────────────────────────
\tcbuselibrary{skins, breakable}

\newtcolorbox{consigna}[1]{
  colback=taskbg,
  colframe=taskborder,
  fonttitle=\bfseries,
  title=#1,
  breakable
}

\newtcolorbox{atencion}{
  colback=alertbg,
  colframe=alertborder,
  fonttitle=\bfseries,
  title={Atención},
  breakable
}

% ── Encabezado ─────────────────────────────────────────────────────────────────
\pagestyle{fancy}
\fancyhf{}
\rhead{Tópico de Especialidad II}
\lhead{Trabajo Práctico — MQTT Avanzado}
\rfoot{Página \thepage}

\hypersetup{colorlinks=true, linkcolor=blue, urlcolor=blue}

% ==============================================================================
\begin{document}
% ==============================================================================

\begin{titlepage}
  \centering
  \vspace*{1.5cm}
  {\large Universidad — Ingeniería en Sistemas\par}
  \vspace{0.3cm}
  {\large Tópico de Especialidad II\par}
  \vspace{1.5cm}
  \rule{\linewidth}{0.5pt}\\[0.4cm]
  {\LARGE\bfseries Trabajo Práctico N.\textsuperscript{o} 2\par}
  \vspace{0.4cm}
  {\Large MQTT Avanzado: QoS, Topics Jerárquicos y Wildcards\par}
  \vspace{0.4cm}
  {\large Extensión del Sistema IoT Agrícola\par}
  \rule{\linewidth}{0.5pt}\\[1.5cm]
  \begin{tabular}{ll}
    \textbf{Modalidad:}   & Grupal (2 a 3 integrantes) \\[4pt]
    \textbf{Entrega:}     & Según cronograma de la cátedra \\[4pt]
    \textbf{Formato:}     & Repositorio + informe PDF + demo en vivo \\
  \end{tabular}
  \vfill
  {\large \today\par}
\end{titlepage}

\tableofcontents
\newpage

% ==============================================================================
\section{Introducción}
% ==============================================================================

En el trabajo práctico anterior se construyó el proyecto \textbf{iot\_agro\_project},
un sistema de sensorización agrícola que combina el protocolo MQTT para la
telemetría de campo con una REST API para la consulta de datos desde un dashboard
Streamlit, persistiendo la información en MongoDB.

Si bien el sistema funciona correctamente en su forma actual, presenta tres
limitaciones técnicas que lo alejan de un escenario de producción real:

\begin{enumerate}
  \item \textbf{Sin garantías de entrega:} los mensajes MQTT se publican con QoS 0
        (\textit{fire and forget}). Si el broker o la red fallan en el momento de
        la transmisión, los datos se pierden sin ningún mecanismo de reenvío.

  \item \textbf{Sin diferenciación de cultivos:} todos los sensores publican en un
        único topic genérico (\texttt{campo/sensores}), sin distinguir a qué zona
        o cultivo pertenece cada lectura. El sistema no puede responder preguntas
        como \textit{``¿cuál es la humedad actual en la zona de tomate?''}.

  \item \textbf{Sin uso de wildcards:} el suscriptor se conecta a un topic fijo y
        exacto. Agregar una nueva zona requiere modificar el código fuente del
        suscriptor, lo cual no es escalable.
\end{enumerate}

El objetivo de este trabajo es que el grupo comprenda, implemente y verifique
las soluciones a estas tres limitaciones utilizando las capacidades nativas
del protocolo MQTT.

% ==============================================================================
\section{Marco Teórico}
% ==============================================================================

Antes de realizar cualquier modificación al código, el grupo debe tener en claro
los conceptos que se explican a continuación. Se recomienda leer esta sección
completa y discutirla en grupo antes de comenzar la implementación.

% ------------------------------------------------------------------------------
\subsection{Quality of Service (QoS)}
% ------------------------------------------------------------------------------

QoS es el mecanismo de MQTT que define el contrato de entrega entre cliente
y broker. Establece cuántos intentos de entrega se realizan y qué confirmaciones
se intercambian en cada caso.

\subsubsection{QoS 0 — At most once}

\begin{lstlisting}[style=texto]
Publisher          Broker          Subscriber
    |                |                 |
    |--PUBLISH------->|                 |
    |   (sin ACK)     |--PUBLISH------->|
    |                 |   (sin ACK)     |
\end{lstlisting}

El mensaje se transmite una única vez sin ninguna confirmación. Si el broker
no está disponible o la red falla durante la transmisión, el mensaje se pierde
permanentemente. Es el nivel por defecto en \texttt{paho-mqtt} cuando no se
especifica el parámetro \texttt{qos}.

\textbf{Consecuencia práctica:} en una red WiFi con interferencias o en un sensor
alimentado por batería que puede perder conexión, un porcentaje de lecturas nunca
llegará al sistema. Esto es aceptable si la frecuencia de muestreo es muy alta
y perder muestras individuales no afecta el análisis.

\subsubsection{QoS 1 — At least once}

\begin{lstlisting}[style=texto]
Publisher          Broker
    |                |
    |--PUBLISH(id=5)->|   El publisher asigna un ID al mensaje
    |<--PUBACK(id=5)--|   El broker confirma recepcion
    |                |
    (si no llega PUBACK, el publisher reenvía el mensaje)
\end{lstlisting}

El publisher almacena el mensaje hasta recibir la confirmación \texttt{PUBACK}
del broker. Si no la recibe en el tiempo esperado, reenvía el mensaje. Esto
garantiza que el mensaje llega \textbf{al menos una vez}, pero puede llegar
\textbf{duplicado} si el \texttt{PUBACK} se pierde en tránsito (el broker ya
recibió el mensaje pero la confirmación no llegó al publisher).

\textbf{Consecuencia práctica:} en el contexto de sensores agrícolas, un mensaje
duplicado significa una lectura duplicada en MongoDB, lo cual puede filtrarse
fácilmente en la consulta. Es el balance razonable entre confiabilidad y overhead.

\subsubsection{QoS 2 — Exactly once}

\begin{lstlisting}[style=texto]
Publisher          Broker
    |                |
    |--PUBLISH------->|   Paso 1: envio
    |<--PUBREC--------|   Paso 2: broker acusa recibo
    |--PUBREL-------->|   Paso 3: publisher libera
    |<--PUBCOMP-------|   Paso 4: broker confirma finalizacion
\end{lstlisting}

Un protocolo de cuatro pasos que garantiza entrega exactamente una vez, sin
pérdidas ni duplicados. El costo es cuatro tramas TCP por mensaje en lugar de
una (QoS 0) o dos (QoS 1).

\textbf{Consecuencia práctica:} apropiado para comandos a actuadores (activar
una válvula de riego, aplicar un fertilizante) donde duplicar la acción tendría
consecuencias reales. Para telemetría de sensores es generalmente excesivo.

\subsubsection{Regla del mínimo}

El nivel de QoS efectivo de una entrega es siempre el \textbf{mínimo} entre
el QoS declarado por el publisher y el QoS declarado por el subscriber. Si el
sensor publica con QoS 1 pero el subscriber se suscribe con QoS 0, la entrega
final al subscriber no tendrá confirmación.

\begin{table}[h]
\centering
\begin{tabular}{lccc}
\toprule
\textbf{Característica}   & \textbf{QoS 0} & \textbf{QoS 1} & \textbf{QoS 2} \\
\midrule
Garantía de entrega       & Ninguna         & Al menos 1 vez  & Exactamente 1 vez \\
Posible pérdida           & Sí              & No              & No \\
Posible duplicado         & No              & Sí              & No \\
Tramas TCP por mensaje    & 1               & 2               & 4 \\
Overhead de red           & Mínimo          & Medio           & Alto \\
\bottomrule
\end{tabular}
\caption{Comparativa de niveles de QoS en MQTT}
\end{table}

% ------------------------------------------------------------------------------
\subsection{Topics Jerárquicos}
% ------------------------------------------------------------------------------

Un topic MQTT es una cadena de texto que define el \textit{canal} por el que
circula un mensaje. El broker lo usa para enrutar cada publicación hacia los
suscriptores correspondientes.

Los topics pueden organizarse en jerarquías separando niveles con el
carácter \textbf{/}:

\begin{lstlisting}[style=texto]
campo/tomate/sensores
  ^      ^       ^
nivel1  nivel2  nivel3
\end{lstlisting}

Esta estructura es arbitraria —el broker no impone ningún esquema— pero seguir
una convención jerárquica tiene ventajas concretas:

\begin{itemize}
  \item Permite filtrar mensajes con precisión usando wildcards (ver sección
        siguiente).
  \item Hace el sistema más legible y autodocumentado.
  \item Facilita el control de acceso: en un broker con autenticación,
        es posible dar permiso a un cliente solo sobre ciertos prefijos.
\end{itemize}

\textbf{Ejemplo aplicado a este proyecto:} en lugar de un topic plano y genérico,
se podría diseñar la siguiente jerarquía:

\begin{lstlisting}[style=texto]
campo/tomate/sensores
campo/zanahoria/sensores
campo/maiz/sensores
\end{lstlisting}

Con este diseño, cada lectura lleva implícita en su topic la zona de cultivo a
la que pertenece.

% ------------------------------------------------------------------------------
\subsection{Wildcards MQTT}
% ------------------------------------------------------------------------------

Los wildcards son caracteres especiales que solo pueden usarse en las
\textbf{suscripciones}, nunca en las publicaciones.

\subsubsection{Wildcard de un nivel: \texttt{+}}

Reemplaza exactamente un nivel en la jerarquía del topic.

\begin{lstlisting}[style=texto]
Patron de suscripcion:   campo/+/sensores

Coincide con:
  campo/tomate/sensores      <- nivel 2 = "tomate"
  campo/zanahoria/sensores   <- nivel 2 = "zanahoria"
  campo/maiz/sensores        <- nivel 2 = "maiz"

NO coincide con:
  campo/sensores             <- solo 2 niveles
  campo/tomate/zona1/sens    <- 4 niveles
\end{lstlisting}

\subsubsection{Wildcard multinivel: \texttt{\#}}

Reemplaza cualquier cantidad de niveles desde el punto donde se coloca.
Solo puede aparecer al final del topic.

\begin{lstlisting}[style=texto]
Patron de suscripcion:   campo/#

Coincide con:
  campo/tomate/sensores
  campo/zanahoria/sensores
  campo/actuadores/riego     <- TAMBIEN capturaría esto
  campo/sistema/alertas      <- y esto
\end{lstlisting}

\subsubsection{Diferencia clave entre \texttt{+} y \texttt{\#}}

El wildcard \texttt{+} es \textbf{selectivo}: captura mensajes que cumplan
exactamente la estructura definida. El wildcard \texttt{\#} es \textbf{amplio}:
captura todo lo que esté bajo un prefijo dado, incluyendo topics que podrían
tener semánticas completamente distintas (sensores, actuadores, alertas, etc.).

En una arquitectura bien diseñada, usar \texttt{+} en lugar de \texttt{\#}
es preferible porque limita el alcance del suscriptor al tipo de dato
que realmente necesita procesar.

% ==============================================================================
\section{Estado Actual del Proyecto}
% ==============================================================================

El código entregado en el TP anterior presenta la siguiente configuración MQTT:

\textbf{En \texttt{sensors/sensor.py}} — el sensor publica en un topic fijo
sin especificar QoS:

\begin{lstlisting}[style=python]
client.publish("campo/sensores", json.dumps(data))
# QoS no especificado -> usa QoS 0 por defecto
# Topic fijo: no hay distincion de zona o cultivo
\end{lstlisting}

\textbf{En \texttt{mqtt\_client/subscriber.py}} — el suscriptor se conecta
a un topic exacto sin wildcard:

\begin{lstlisting}[style=python]
client.subscribe("campo/sensores")
# Topic exacto: no usa wildcard
# Si se agrega una nueva zona, hay que modificar este archivo
\end{lstlisting}

\textbf{En \texttt{docker-compose.yml}} — los sensores solo reciben
\texttt{SENSOR\_ID} como variable de entorno, sin ninguna referencia a zona:

\begin{lstlisting}[style=texto]
sensor1:
  environment:
    - SENSOR_ID=1

sensor2:
  environment:
    - SENSOR_ID=2
\end{lstlisting}

\begin{atencion}
El grupo \textbf{no debe modificar} ningún archivo del proyecto base antes de
haber leído y comprendido completamente la sección de consignas. Los cambios
deben realizarse de forma ordenada, justificando cada decisión en el informe.
\end{atencion}

% ==============================================================================
\section{Consignas}
% ==============================================================================

% ------------------------------------------------------------------------------
\subsection{Parte 1 — Implementar QoS 1 en el sistema}
% ------------------------------------------------------------------------------

\begin{consigna}{Consigna 1.1 — Activar QoS 1 en la publicación}
Modificar \texttt{sensors/sensor.py} para que cada mensaje se publique con
QoS nivel 1. La API de \texttt{paho-mqtt} recibe el nivel de QoS como parámetro
del método \texttt{publish()}. Investigar la firma del método y aplicar el
cambio correspondiente.
\end{consigna}

\begin{consigna}{Consigna 1.2 — Activar QoS 1 en la suscripción}
Modificar \texttt{mqtt\_client/subscriber.py} para que la suscripción al
topic se realice con QoS nivel 1. El método \texttt{subscribe()} de
\texttt{paho-mqtt} acepta el nivel de QoS como argumento. Investigar y aplicar.
\end{consigna}

\begin{consigna}{Consigna 1.3 — Verificar el QoS recibido}
El objeto \texttt{msg} que llega al callback \texttt{on\_message()} contiene
el atributo \texttt{msg.qos} que indica el QoS con el que fue entregado ese
mensaje. Modificar el callback para que:
\begin{itemize}
  \item Imprima por consola el QoS de cada mensaje recibido.
  \item Almacene el valor de \texttt{msg.qos} como campo en el documento
        de MongoDB junto con los demás datos del sensor.
\end{itemize}
Verificar en MongoDB Compass que el campo \texttt{qos} aparece en los
documentos con el valor esperado.
\end{consigna}

\textbf{Preguntas a responder en el informe:}
\begin{itemize}
  \item ¿Por qué se eligió QoS 1 y no QoS 2 para este sistema?
  \item ¿Qué pasaría si el publisher usa QoS 1 pero el subscriber
        se suscribe con QoS 0? ¿Qué nivel efectivo tendría la entrega?
  \item ¿En qué escenario concreto del dominio agrícola sería necesario
        usar QoS 2?
\end{itemize}

% ------------------------------------------------------------------------------
\subsection{Parte 2 — Topics jerárquicos por zona de cultivo}
% ------------------------------------------------------------------------------

\begin{consigna}{Consigna 2.1 — Agregar la variable de entorno ZONA}
Modificar \texttt{docker-compose.yml} para que cada servicio sensor reciba,
además de \texttt{SENSOR\_ID}, una nueva variable de entorno llamada
\texttt{ZONA}. Definir al menos tres zonas distintas, por ejemplo:
\texttt{tomate}, \texttt{zanahoria} y \texttt{maiz}. Puede renombrar
los servicios para que el nombre refleje la zona correspondiente.
\end{consigna}

\begin{consigna}{Consigna 2.2 — Construir el topic dinámicamente}
Modificar \texttt{sensors/sensor.py} para que:
\begin{itemize}
  \item Lea la variable de entorno \texttt{ZONA} al iniciar.
  \item Construya el topic de publicación usando la estructura
        \texttt{campo/<zona>/sensores}, donde \texttt{<zona>} proviene
        de la variable de entorno.
  \item Incluya el valor de \texttt{zona} en el payload JSON que publica
        para que el dato viaje también en el mensaje.
\end{itemize}
El mismo archivo \texttt{sensor.py} debe poder ser reutilizado por todos
los contenedores de sensores sin modificación. Solo las variables de entorno
deben diferenciar su comportamiento.
\end{consigna}

\begin{consigna}{Consigna 2.3 — Verificar los topics en el broker}
Con el sistema levantado, utilizar un cliente MQTT externo compatible con
\textbf{AWS IoT Core} (por ejemplo, \textbf{MQTT Explorer} con TLS y los certificados del thing,
o la \textbf{Consola $\rightarrow$ MQTT test client} en AWS IoT Core) para suscribirse
a cada topic individualmente y verificar que los mensajes llegan al topic correcto.
\textit{No aplica} suscripción anónima a \texttt{localhost:1883}: la consigna del Hito 2 exige el broker
gestionado en la nube, no Mosquitto local.

Ejemplo de forma genérica (endpoint y rutas de certificado según el entorno):
\begin{lstlisting}[style=texto]
# Cliente con TLS al endpoint a2apsmaa0mdv52-ats.iot.*.amazonaws.com:8883
# (no usar broker local :1883 para la entrega Hito 2)
\end{lstlisting}
\end{consigna}

\textbf{Preguntas a responder en el informe:}
\begin{itemize}
  \item ¿Qué ventaja tiene diseñar la jerarquía de topics como
        \texttt{campo/<zona>/sensores} en lugar de usar un topic
        plano como \texttt{sensores-campo}?
  \item ¿Cómo afecta este cambio al subscriber? ¿Sigue funcionando
        sin modificaciones? Justificar.
\end{itemize}

% ------------------------------------------------------------------------------
\subsection{Parte 3 — Wildcard en el suscriptor}
% ------------------------------------------------------------------------------

\begin{consigna}{Consigna 3.1 — Suscripción con wildcard \texttt{+}}
Modificar \texttt{mqtt\_client/subscriber.py} para que, en lugar de
suscribirse a un topic fijo, utilice el wildcard \texttt{+} para
capturar mensajes de todas las zonas con un único \texttt{subscribe()}.

El patrón de suscripción debe tener la forma \texttt{campo/+/sensores}.
\end{consigna}

\begin{consigna}{Consigna 3.2 — Extraer la zona desde el topic recibido}
Cuando el broker entrega un mensaje al callback \texttt{on\_message()},
el objeto \texttt{msg} tiene el atributo \texttt{msg.topic} que contiene
el topic exacto por el que llegó ese mensaje (por ejemplo,
\texttt{campo/tomate/sensores}).

Modificar el callback para que:
\begin{itemize}
  \item Extraiga el nombre de la zona del topic recibido usando manipulación
        de strings (\texttt{split()}).
  \item Almacene la zona como campo \texttt{zona} en el documento de MongoDB,
        independientemente de si el payload ya la incluye o no.
\end{itemize}
\end{consigna}

\begin{consigna}{Consigna 3.3 — Demostrar la escalabilidad}
Sin modificar el código del subscriber, agregar un cuarto sensor con una
nueva zona (por ejemplo, \texttt{lechuga}) solo en \texttt{docker-compose.yml}
y verificar que:
\begin{itemize}
  \item El subscriber captura automáticamente los mensajes de la nueva zona.
  \item Los documentos de esa zona aparecen correctamente en MongoDB.
\end{itemize}
Esto demuestra que el diseño con wildcard es escalable.
Documentar este experimento en el informe con capturas de logs.
\end{consigna}

\textbf{Preguntas a responder en el informe:}
\begin{itemize}
  \item ¿Por qué se usa \texttt{campo/+/sensores} en lugar de \texttt{campo/\#}?
        Describir al menos un escenario donde \texttt{campo/\#} capturaría
        mensajes no deseados.
  \item ¿Cómo debería modificarse el subscriber si en el futuro se quisiera
        ignorar mensajes de una zona específica?
\end{itemize}

% ------------------------------------------------------------------------------
\subsection{Parte 4 — Extensión de la REST API}
% ------------------------------------------------------------------------------

\begin{consigna}{Consigna 4.1 — Endpoint de filtrado por zona}
La REST API actual solo expone \texttt{GET /logs} que devuelve las últimas
20 lecturas sin importar la zona. Agregar en \texttt{rest\_api/app.py}
los siguientes endpoints:
\begin{itemize}
  \item \texttt{GET /logs/<zona>} — devuelve las últimas 20 lecturas de
        una zona específica, filtrando por el campo \texttt{zona} en MongoDB.
  \item \texttt{GET /zonas} — devuelve la lista de zonas distintas presentes
        en la colección, usando la operación \texttt{distinct()} de PyMongo.
\end{itemize}
Verificar ambos endpoints con \texttt{curl} o un cliente REST (Postman,
Insomnia, etc.) e incluir las respuestas en el informe.
\end{consigna}

% ------------------------------------------------------------------------------
\subsection{Parte 5 — Dashboard Streamlit separado por zonas}
% ------------------------------------------------------------------------------

\begin{consigna}{Consigna 5.1 — Tabs por zona de cultivo}
Modificar \texttt{frontend/app.py} para reemplazar la vista única por un
dashboard organizado en tabs, donde cada tab corresponde a una zona de cultivo.
Los requerimientos son:

\begin{enumerate}
  \item Las zonas disponibles deben obtenerse dinámicamente consultando
        el endpoint \texttt{GET /zonas} de la API, no estar hardcodeadas
        en el código del frontend.
  \item Cada tab debe mostrar únicamente las lecturas de su zona correspondiente,
        consultando \texttt{GET /logs/<zona>}.
  \item Cada tab debe incluir al menos:
        \begin{itemize}
          \item Una métrica con la temperatura promedio de las últimas lecturas.
          \item Una métrica con la humedad promedio.
          \item Una tabla con las lecturas individuales.
          \item Un gráfico de línea con la evolución de temperatura y humedad.
        \end{itemize}
  \item Debe existir un tab adicional llamado ``Todas las zonas'' que muestre
        las últimas 20 lecturas globales (\texttt{GET /logs}) con un gráfico
        comparativo por zona.
\end{enumerate}

Si se agrega una nueva zona al sistema (nueva entrada en \texttt{docker-compose.yml}),
el dashboard debe mostrar el nuevo tab automáticamente sin requerir cambios en
el código del frontend.
\end{consigna}

% ==============================================================================
\section{Criterios de Evaluación}
% ==============================================================================

\begin{table}[h]
\centering
\begin{tabular}{p{7cm}c}
\toprule
\textbf{Criterio} & \textbf{Puntaje} \\
\midrule
Parte 1: QoS 1 implementado y verificado en MongoDB & 15 pts \\
Parte 2: Topics por zona funcionando correctamente & 20 pts \\
Parte 3: Wildcard y escalabilidad demostrada & 20 pts \\
Parte 4: Endpoints REST por zona & 15 pts \\
Parte 5: Dashboard Streamlit por zonas & 20 pts \\
Calidad del informe y respuestas a preguntas & 10 pts \\
\midrule
\textbf{Total} & \textbf{100 pts} \\
\bottomrule
\end{tabular}
\caption{Distribución de puntaje}
\end{table}

% ==============================================================================
\section{Formato de Entrega}
% ==============================================================================

\subsection{Repositorio}

El grupo debe entregar un repositorio (GitHub, GitLab o similar) con:
\begin{itemize}
  \item El código fuente completo del proyecto con todos los cambios aplicados.
  \item Un archivo \texttt{README.md} actualizado con las instrucciones para
        levantar el sistema.
  \item Un archivo \texttt{.env.example} o similar documentando todas las
        variables de entorno utilizadas.
\end{itemize}

\subsection{Informe}

Un documento PDF que incluya:
\begin{enumerate}
  \item Portada con integrantes del grupo y fecha de entrega.
  \item Para cada parte: descripción de los cambios realizados, fragmentos
        de código relevantes y capturas de pantalla o logs que demuestren
        el funcionamiento.
  \item Respuestas fundamentadas a todas las preguntas planteadas en las
        consignas.
  \item Una sección de conclusiones grupales sobre las ventajas y desventajas
        de los mecanismos implementados (QoS, topics jerárquicos, wildcards)
        en el contexto de sistemas IoT reales.
\end{enumerate}

\subsection{Demo en vivo}

Durante la instancia de defensa, el grupo debe poder:
\begin{itemize}
  \item Levantar el sistema completo con \texttt{docker compose up --build}.
  \item Mostrar en tiempo real los logs de los sensores y el subscriber.
  \item Navegar el dashboard Streamlit mostrando los tabs por zona.
  \item Agregar un nuevo sensor con una zona nueva y demostrar que el
        subscriber y el dashboard lo incorporan automáticamente.
  \item Mostrar los documentos almacenados en MongoDB Compass con los
        campos \texttt{zona} y \texttt{qos}.
\end{itemize}

% ==============================================================================
\section{Recursos de Consulta}
% ==============================================================================

\begin{itemize}
  \item OASIS MQTT Technical Committee. \textit{MQTT Version 5.0 Specification}.
        \url{https://docs.oasis-open.org/mqtt/mqtt/v5.0/mqtt-v5.0.html}
  \item Eclipse Paho. \textit{paho-mqtt Python Client}.
        \url{https://eclipse.dev/paho/files/paho.mqtt.python/html/index.html}
  \item HiveMQ. \textit{MQTT Essentials — QoS Levels}.
        \url{https://www.hivemq.com/mqtt-essentials/}
  \item MongoDB Inc. \textit{PyMongo — Query Documents}.
        \url{https://pymongo.readthedocs.io/en/stable/tutorial.html}
  \item Streamlit Inc. \textit{st.tabs — API reference}.
        \url{https://docs.streamlit.io/develop/api-reference/layout/st.tabs}
\end{itemize}

\begin{atencion}
Queda prohibido el uso de herramientas de generación automática de código
(LLMs, Copilot, etc.) para resolver las consignas. El informe debe demostrar
comprensión genuina de cada cambio realizado. Durante la defensa se podrán
realizar preguntas individuales a cualquier integrante del grupo.
\end{atencion}

\end{document}
